\centerline{\bf \Huge Разнобой}

\begin{flushright}
        \scriptsize{с новой морозой вас}
\end{flushright}

\problem{1}
\textbf{Считаем в синусах!}
В треугольнике $A B C$ с углами $\angle A=\alpha, \angle B=\beta, \angle C=\gamma$ и радиусом описанной окружности $R$ провели все а) высоты; б) биссектрисы. Найдите длины всех отрезков на рисунке.

\problem{2}
\textbf{Я сказал считаем в синусах!}
В остроугольном треугольнике $A B C$ проведены высоты $A A^{\prime}$ и $B B^{\prime}$, точка $O$ - центр его описанной окружности. Докажите, что расстояние от точки $A^{\prime}$ до прямой $B O$ равно расстоянию от точки $B^{\prime}$ до прямой $A O$.

\problem{3}
Группа ЭлМШ из 20 учеников пришла писать муниципальный этап. Организаторы что-то напутали и задач в варианте оказалось 6. Верно ли следующее утверждение: по результатам олимпиады всегда найдутся либо какие-то пять учеников и такие две задачи, что каждый из этих пять учеников решил эти две задачи, либо же всегда найдутся какие-то пять учеников и такие две задачи, что каждый из этих 5 учеников не решил ни одну из этих двух задач.
%из патнема


\problem{4}
Докажите, что стороны любого неравнобедренного треугольника можно либо все увеличить, либо все уменьшить на одну и ту же величину так, чтобы получился прямоугольный треугольник.

%https://math.mosolymp.ru/upload/files/2020/khamovniki/9/9-1/2019-09-09_Resheniya.pdf

\problem{5}
Алина написала муницип из 5 задач, где каждая задача оценивалась от 0 до 7 баллов, но она не хочет говорить свои баллы Расулу Владимировичу. Про каждую задачу он может задать вопрос ответ на который да или нет. Алина поклялась отвечать правду. За какое наименьшее количество вопросов РВ сможет гарантированно отгадать баллы Алины?
%https://math.mosolymp.ru/upload/files/2020/khamovniki/9/9-1/2020-03-23_Raznoboy__Resheniya.pdf

\problem{6}
Дан параллелограмм $A B C D$ с острым углом $A$. На дуге $A B$ описанной окружности треугольника $A B D$, не содержащей точку $D$, выбрана произвольная точка $X$. Лучи $X A$ и $C D$ пересекаются в точке $Y$. Докажите, что центр описанной окружности треугольника $C X Y$ лежит на прямой, проходящей через $D$ перпендикулярно $A D$.

%https://math.mosolymp.ru/upload/files/2021/khamovniki/10/10-2/2021-05-21_Otborochnaya_olimpiada,_resheniya.pdf